\chapter{Taylor Expansions, Asymptotic Developments \& Riemann Integration}
\label{chap:exo_ch06}

\begin{quicksummary}{Leibniz Formula, Taylor-Young, Riemann Sums \& Integration Techniques}
\begin{enumerate}[leftmargin=1.4em,itemsep=3pt]
  \item \textbf{Leibniz Formula for $n$-th Derivative:}
  \[
  (u \cdot v)^{(n)} = \sum_{k=0}^n \binom{n}{k} u^{(k)} v^{(n-k)}.
  \]
  \item \textbf{Taylor-Young Formula at $0$ (Maclaurin DL$_n(0)$):}
  \[
  f(x) = \sum_{k=0}^n \frac{f^{(k)}(0)}{k!} x^k + o(x^n) \quad \text{as } x \to 0.
  \]
  \item \textbf{Reference Taylor Expansions at $0$:}
  \[
  e^x = 1 + x + \frac{x^2}{2!} + \frac{x^3}{3!} + \dots + \frac{x^n}{n!} + o(x^n), \qquad \cos(x) = 1 - \frac{x^2}{2!} + \frac{x^4}{4!} - \dots + o(x^{2n+1}).
  \]
  \[
  \sin(x) = x - \frac{x^3}{3!} + \frac{x^5}{5!} - \dots + o(x^{2n+2}), \qquad \ln(1 + x) = x - \frac{x^2}{2} + \frac{x^3}{3} - \dots + (-1)^{n-1}\frac{x^n}{n} + o(x^n).
  \]
  \[
  \frac{1}{1 - x} = 1 + x + x^2 + \dots + x^n + o(x^n), \qquad (1 + x)^\alpha = 1 + \alpha x + \frac{\alpha(\alpha-1)}{2} x^2 + \dots + o(x^n).
  \]
  \item \textbf{Riemann Sum Limit Theorem:}
  If $f$ is continuous on $[a, b]$:
  \[
  \lim_{n \to \infty} \frac{b - a}{n} \sum_{k=1}^n f\left(a + k\frac{b - a}{n}\right) = \int_a^b f(x)\,\dd x.
  \]
  \item \textbf{Fundamental Theorem of Calculus:}
  If $f \in C([a, b])$ and $F(x) = \int_a^x f(t)\,\dd t$, then $F \in C^1([a, b])$ and $F'(x) = f(x)$.
  \item \textbf{Integration by Parts (IPP) \& Substitution:}
  \[
  \int u(x) v'(x)\,\dd x = u(x) v(x) - \int u'(x) v(x)\,\dd x, \qquad \int_a^b f(\phi(t))\phi'(t)\,\dd t = \int_{\phi(a)}^{\phi(b)} f(u)\,\dd u.
  \]
\end{enumerate}
\end{quicksummary}

\vspace{10pt}

% ===========================================================================
% EXERCISE 6.1
% ===========================================================================
\begin{exercise}[$\bigstar$ Foundational]{The Leibniz Formula for the $n$-th Derivative}
\begin{enumerate}[leftmargin=1.5em,itemsep=2pt]
  \item State the Leibniz formula for the $n$-th derivative of a product of two $n$-times differentiable functions $u$ and $v$.
  \item Calculate the $n$-th derivative $f^{(n)}(x)$ for $f(x) = x^2 e^{3x}$ ($n \ge 2$).
  \item Calculate the $n$-th derivative $g^{(n)}(x)$ for $g(x) = (x^2 + 1) \cos(x)$.
\end{enumerate}
\end{exercise}

\begin{tabletworkspace}[5.0cm]
\end{tabletworkspace}

\begin{solution}
\textbf{1. Leibniz Formula:}
If $u, v \in C^n(I)$:
\[
(uv)^{(n)}(x) = \sum_{k=0}^n \binom{n}{k} u^{(k)}(x) v^{(n-k)}(x).
\]

\textbf{2. $n$-th Derivative of $f(x) = x^2 e^{3x}$:}
Let $u(x) = x^2$ and $v(x) = e^{3x}$.
Compute the successive derivatives of $u(x)$:
\[
u^{(0)}(x) = x^2, \quad u'(x) = 2x, \quad u''(x) = 2, \quad \text{and } u^{(k)}(x) = 0 \quad \text{for all } k \ge 3.
\]
The derivatives of $v(x) = e^{3x}$ are:
\[
v^{(m)}(x) = 3^m e^{3x}.
\]
Applying the Leibniz formula (since all terms for $k \ge 3$ vanish):
\begin{align*}
f^{(n)}(x) &= \binom{n}{0} u(x) v^{(n)}(x) + \binom{n}{1} u'(x) v^{(n-1)}(x) + \binom{n}{2} u''(x) v^{(n-2)}(x) \\
&= 1 \cdot x^2 \cdot 3^n e^{3x} + n \cdot (2x) \cdot 3^{n-1} e^{3x} + \frac{n(n-1)}{2} \cdot 2 \cdot 3^{n-2} e^{3x} \\
&= 3^{n-2} e^{3x} \left[ 9x^2 + 6nx + n(n-1) \right]. \quad \blacksquare
\end{align*}

\textbf{3. $n$-th Derivative of $g(x) = (x^2 + 1)\cos(x)$:}
Let $u(x) = x^2 + 1$ and $v(x) = \cos(x)$.
Derivatives of $u$: $u'(x) = 2x$, $u''(x) = 2$, $u^{(k)}(x) = 0$ for $k \ge 3$.
Derivatives of $v$: $v^{(m)}(x) = \cos\left(x + m\frac{\pi}{2}\right)$.
Applying Leibniz:
\begin{align*}
g^{(n)}(x) &= \binom{n}{0} (x^2 + 1) \cos\left(x + \frac{n\pi}{2}\right) + \binom{n}{1} (2x) \cos\left(x + \frac{(n-1)\pi}{2}\right) \\
&\quad + \binom{n}{2} (2) \cos\left(x + \frac{(n-2)\pi}{2}\right) \\
&= (x^2 + 1)\cos\left(x + \frac{n\pi}{2}\right) + 2nx \sin\left(x + \frac{n\pi}{2}\right) - n(n-1)\cos\left(x + \frac{n\pi}{2}\right) \\
&= \left[ x^2 + 1 - n(n-1) \right] \cos\left(x + \frac{n\pi}{2}\right) + 2nx \sin\left(x + \frac{n\pi}{2}\right). \quad \blacksquare
\end{align*}
\end{solution}

% ===========================================================================
% EXERCISE 6.2
% ===========================================================================
\begin{exercise}[$\bigstar$ Foundational]{Direct Product and Composition of Maclaurin Series}
Compute the Taylor-Young asymptotic expansion (development) at order 4 at $x_0 = 0$ for:
\begin{enumerate}[leftmargin=1.5em,itemsep=2pt]
  \item $f(x) = e^x \cos(x)$ at order 4.
  \item $g(x) = \ln(1 + \sin(x))$ at order 3.
  \item $h(x) = \frac{1}{\sqrt{1 - x^2}}$ at order 4.
\end{enumerate}
\end{exercise}

\begin{tabletworkspace}[5.0cm]
\end{tabletworkspace}

\begin{solution}
\textbf{1. Expansion of $f(x) = e^x \cos(x)$ to Order 4:}
The reference expansions at $0$ are:
\[
e^x = 1 + x + \frac{x^2}{2} + \frac{x^3}{6} + \frac{x^4}{24} + o(x^4),
\]
\[
\cos(x) = 1 - \frac{x^2}{2} + \frac{x^4}{24} + o(x^4).
\]
Multiplying and retaining terms of order $\le 4$:
\begin{align*}
e^x \cos(x) &= \left( 1 + x + \frac{x^2}{2} + \frac{x^3}{6} + \frac{x^4}{24} \right) \left( 1 - \frac{x^2}{2} + \frac{x^4}{24} \right) + o(x^4) \\
&= 1\left(1 - \frac{x^2}{2} + \frac{x^4}{24}\right) + x\left(1 - \frac{x^2}{2}\right) + \frac{x^2}{2}\left(1 - \frac{x^2}{2}\right) + \frac{x^3}{6}(1) + \frac{x^4}{24}(1) + o(x^4) \\
&= 1 - \frac{x^2}{2} + \frac{x^4}{24} + x - \frac{x^3}{2} + \frac{x^2}{2} - \frac{x^4}{4} + \frac{x^3}{6} + \frac{x^4}{24} + o(x^4) \\
&= 1 + x + \left(-\frac{1}{2} + \frac{1}{2}\right)x^2 + \left(-\frac{1}{2} + \frac{1}{6}\right)x^3 + \left(\frac{1}{24} - \frac{1}{4} + \frac{1}{24}\right)x^4 + o(x^4) \\
&= 1 + x - \frac{x^3}{3} - \frac{x^4}{6} + o(x^4). \quad \blacksquare
\end{align*}

\textbf{2. Expansion of $g(x) = \ln(1 + \sin x)$ to Order 3:}
We know:
\[
u = \sin(x) = x - \frac{x^3}{6} + o(x^3) \xrightarrow{x \to 0} 0.
\]
Expansion of $\ln(1 + u) = u - \frac{u^2}{2} + \frac{u^3}{3} + o(u^3)$:
Compute powers of $u$:
\[
u = x - \frac{x^3}{6} + o(x^3), \qquad u^2 = \left(x - \frac{x^3}{6}\right)^2 = x^2 + o(x^3), \qquad u^3 = x^3 + o(x^3).
\]
Substitute:
\begin{align*}
\ln(1 + \sin x) &= \left(x - \frac{x^3}{6}\right) - \frac{1}{2}\left(x^2\right) + \frac{1}{3}\left(x^3\right) + o(x^3) \\
&= x - \frac{x^2}{2} + \left(-\frac{1}{6} + \frac{1}{3}\right)x^3 + o(x^3) = x - \frac{x^2}{2} + \frac{x^3}{6} + o(x^3). \quad \blacksquare
\end{align*}

\textbf{3. Expansion of $h(x) = (1 - x^2)^{-1/2}$ to Order 4:}
Using $(1 + u)^\alpha = 1 + \alpha u + \frac{\alpha(\alpha - 1)}{2} u^2 + o(u^2)$ with $\alpha = -1/2$ and $u = -x^2$:
\[
\alpha = -\frac{1}{2}, \qquad \frac{\alpha(\alpha - 1)}{2} = \frac{(-1/2)(-3/2)}{2} = \frac{3}{8}.
\]
Substituting $u = -x^2$:
\[
(1 - x^2)^{-1/2} = 1 + \left(-\frac{1}{2}\right)(-x^2) + \frac{3}{8}(-x^2)^2 + o(x^4) = 1 + \frac{x^2}{2} + \frac{3x^4}{8} + o(x^4). \quad \blacksquare
\]
\end{solution}

% ===========================================================================
% EXERCISE 6.3
% ===========================================================================
\begin{exercise}[$\bigstar$ Foundational]{Division and Inversion of Asymptotic Expansions}
\begin{enumerate}[leftmargin=1.5em,itemsep=2pt]
  \item Compute the asymptotic expansion of $\tan(x) = \frac{\sin(x)}{\cos(x)}$ to order 5 at $x_0 = 0$ using polynomial long division.
  \item Compute the asymptotic expansion of $f(x) = \frac{1}{1 + e^x}$ to order 3 at $0$.
\end{enumerate}
\end{exercise}

\begin{tabletworkspace}[5.0cm]
\end{tabletworkspace}

\begin{solution}
\textbf{1. Expansion of $\tan(x)$ to Order 5:}
We have $\sin(x) = x - \frac{x^3}{6} + \frac{x^5}{120} + o(x^5)$ and $\cos(x) = 1 - \frac{x^2}{2} + \frac{x^4}{24} + o(x^5)$.
Perform division according to increasing powers:
\[
\begin{array}{r@{\quad}l}
x - \frac{x^3}{6} + \frac{x^5}{120} & \vline\enspace 1 - \frac{x^2}{2} + \frac{x^4}{24} \\
\cline{2-2}
-(x - \frac{x^3}{2} + \frac{x^5}{24}) & x + \frac{x^3}{3} + \frac{2x^5}{15} \\
\cline{1-1}
\frac{x^3}{3} - \frac{x^5}{30} & \\
-(\frac{x^3}{3} - \frac{x^5}{6}) & \\
\cline{1-1}
\frac{2x^5}{15} &
\end{array}
\]
Thus:
\[
\tan(x) = x + \frac{x^3}{3} + \frac{2x^5}{15} + o(x^5). \quad \blacksquare
\]

\textbf{2. Expansion of $f(x) = \frac{1}{1 + e^x}$ to Order 3:}
Expand $e^x = 1 + x + \frac{x^2}{2} + \frac{x^3}{6} + o(x^3)$.
Denominator:
\[
1 + e^x = 2 + x + \frac{x^2}{2} + \frac{x^3}{6} + o(x^3) = 2 \left[ 1 + \left( \frac{x}{2} + \frac{x^2}{4} + \frac{x^3}{12} \right) \right].
\]
Let $u = \frac{x}{2} + \frac{x^2}{4} + \frac{x^3}{12} + o(x^3) \to 0$.
Use $\frac{1}{1 + u} = 1 - u + u^2 - u^3 + o(u^3)$:
Compute powers of $u$:
\[
u = \frac{x}{2} + \frac{x^2}{4} + \frac{x^3}{12} + o(x^3), \qquad u^2 = \frac{x^2}{4} + \frac{x^3}{4} + o(x^3), \qquad u^3 = \frac{x^3}{8} + o(x^3).
\]
Substitute:
\begin{align*}
\frac{1}{1 + u} &= 1 - \left(\frac{x}{2} + \frac{x^2}{4} + \frac{x^3}{12}\right) + \left(\frac{x^2}{4} + \frac{x^3}{4}\right) - \frac{x^3}{8} + o(x^3) \\
&= 1 - \frac{x}{2} + \left(-\frac{1}{4} + \frac{1}{4}\right)x^2 + \left(-\frac{1}{12} + \frac{1}{4} - \frac{1}{8}\right)x^3 + o(x^3) \\
&= 1 - \frac{x}{2} + 0x^2 + \left(\frac{-2 + 6 - 3}{24}\right)x^3 + o(x^3) = 1 - \frac{x}{2} + \frac{x^3}{24} + o(x^3).
\end{align*}
Multiply by $\frac{1}{2}$:
\[
f(x) = \frac{1}{2} - \frac{x}{4} + \frac{x^3}{48} + o(x^3). \quad \blacksquare
\]
\end{solution}

% ===========================================================================
% EXERCISE 6.4
% ===========================================================================
\begin{exercise}[$\bigstar\bigstar$ Intermediate]{Challenging Limits via Taylor Expansions}
Calculate the following indeterminate limits using asymptotic developments:
\begin{enumerate}[leftmargin=1.5em,itemsep=2pt]
  \item $L_1 = \lim_{x \to 0} \frac{\cos(x) - e^{-x^2/2}}{x^4}$.
  \item $L_2 = \lim_{x \to 0} \frac{\tan(x) - \sin(x)}{x^3}$.
  \item $L_3 = \lim_{x \to 0} \left( \frac{1}{x^2} - \frac{1}{\sin^2(x)} \right)$.
\end{enumerate}
\end{exercise}

\begin{tabletworkspace}[5.0cm]
\end{tabletworkspace}

\begin{solution}
\textbf{1. Computation of $L_1$:}
Expand to order 4:
\[
\cos(x) = 1 - \frac{x^2}{2} + \frac{x^4}{24} + o(x^4).
\]
For $e^{-x^2/2}$, let $u = -x^2/2$:
\[
e^u = 1 + u + \frac{u^2}{2} + o(u^2) = 1 + \left(-\frac{x^2}{2}\right) + \frac{1}{2}\left(-\frac{x^2}{2}\right)^2 + o(x^4) = 1 - \frac{x^2}{2} + \frac{x^4}{8} + o(x^4).
\]
Subtract the expansions:
\begin{align*}
\cos(x) - e^{-x^2/2} &= \left(1 - \frac{x^2}{2} + \frac{x^4}{24}\right) - \left(1 - \frac{x^2}{2} + \frac{x^4}{8}\right) + o(x^4) \\
&= \left(\frac{1}{24} - \frac{1}{8}\right)x^4 + o(x^4) = \left(\frac{1 - 3}{24}\right)x^4 + o(x^4) = -\frac{1}{12}x^4 + o(x^4).
\end{align*}
Divide by $x^4$:
\[
L_1 = \lim_{x \to 0} \frac{-\frac{1}{12}x^4 + o(x^4)}{x^4} = -\frac{1}{12}. \quad \blacksquare
\]

\textbf{2. Computation of $L_2$:}
Expand $\tan(x)$ and $\sin(x)$ to order 3:
\[
\tan(x) = x + \frac{x^3}{3} + o(x^3), \qquad \sin(x) = x - \frac{x^3}{6} + o(x^3).
\]
Subtract:
\[
\tan(x) - \sin(x) = \left(x + \frac{x^3}{3}\right) - \left(x - \frac{x^3}{6}\right) + o(x^3) = \left(\frac{1}{3} + \frac{1}{6}\right)x^3 + o(x^3) = \frac{x^3}{2} + o(x^3).
\]
Divide by $x^3$:
\[
L_2 = \lim_{x \to 0} \frac{x^3/2 + o(x^3)}{x^3} = \frac{1}{2}. \quad \blacksquare
\]

\textbf{3. Computation of $L_3$:}
Bring to a common denominator:
\[
\frac{1}{x^2} - \frac{1}{\sin^2 x} = \frac{\sin^2 x - x^2}{x^2 \sin^2 x}.
\]
Denominator: since $\sin x \sim x$, $x^2 \sin^2 x \sim x^4$.
Numerator: expand $\sin(x) = x - \frac{x^3}{6} + o(x^3)$:
\[
\sin^2(x) = \left(x - \frac{x^3}{6} + o(x^3)\right)^2 = x^2 - \frac{x^4}{3} + o(x^4).
\]
Thus:
\[
\sin^2(x) - x^2 = -\frac{x^4}{3} + o(x^4).
\]
Divide by denominator equivalent:
\[
L_3 = \lim_{x \to 0} \frac{-x^4/3 + o(x^4)}{x^4} = -\frac{1}{3}. \quad \blacksquare
\]
\end{solution}

% ===========================================================================
% EXERCISE 6.5
% ===========================================================================
\begin{exercise}[$\bigstar\bigstar$ Intermediate]{Asymptotes and Relative Position of Curves}
Consider the function:
\[
f(x) = \sqrt{x^2 + 2x + 3}.
\]
\begin{enumerate}[leftmargin=1.5em,itemsep=2pt]
  \item Determine the asymptotic expansion of $f(x)$ as $x \to +\infty$ in powers of $\frac{1}{x}$ up to order $\frac{1}{x}$.
  \item Deduce the equation of the oblique asymptote $(\Delta)$ to the curve $\mathcal{C}_f$ at $+\infty$.
  \item Determine the position of the curve $\mathcal{C}_f$ relative to $(\Delta)$ in the neighborhood of $+\infty$.
  \item Determine the asymptote and relative position as $x \to -\infty$.
\end{enumerate}
\end{exercise}

\begin{tabletworkspace}[5.0cm]
\end{tabletworkspace}

\begin{solution}
\textbf{1. Expansion as $x \to +\infty$:}
For $x > 0$, factor $x = \sqrt{x^2}$:
\[
f(x) = x \sqrt{1 + \frac{2}{x} + \frac{3}{x^2}}.
\]
Let $u = \frac{2}{x} + \frac{3}{x^2} \to 0$ as $x \to +\infty$.
Expand $(1 + u)^{1/2} = 1 + \frac{1}{2}u - \frac{1}{8}u^2 + o(u^2)$:
\[
u = \frac{2}{x} + \frac{3}{x^2}, \qquad u^2 = \left(\frac{2}{x} + \frac{3}{x^2}\right)^2 = \frac{4}{x^2} + o\left(\frac{1}{x^2}\right).
\]
Substitute:
\begin{align*}
\sqrt{1 + u} &= 1 + \frac{1}{2}\left(\frac{2}{x} + \frac{3}{x^2}\right) - \frac{1}{8}\left(\frac{4}{x^2}\right) + o\left(\frac{1}{x^2}\right) \\
&= 1 + \frac{1}{x} + \frac{3}{2x^2} - \frac{1}{2x^2} + o\left(\frac{1}{x^2}\right) = 1 + \frac{1}{x} + \frac{1}{x^2} + o\left(\frac{1}{x^2}\right).
\end{align*}
Multiply by $x$:
\[
f(x) = x + 1 + \frac{1}{x} + o\left(\frac{1}{x}\right).
\]

\textbf{2. Oblique Asymptote at $+\infty$:}
The linear part gives the equation of the oblique asymptote:
\[
(\Delta): \quad y = x + 1.
\]

\textbf{3. Relative Position at $+\infty$:}
Compute the vertical distance $f(x) - y_{(\Delta)}$:
\[
f(x) - (x + 1) = \frac{1}{x} + o\left(\frac{1}{x}\right).
\]
As $x \to +\infty$, $\frac{1}{x} > 0$.
Therefore, $f(x) - y_{(\Delta)} > 0$, meaning the curve $\mathcal{C}_f$ lies strictly \textbf{above} the asymptote $(\Delta)$.

\textbf{4. Behavior as $x \to -\infty$:}
For $x < 0$, $\sqrt{x^2} = -x = |x|$.
\[
f(x) = -x \sqrt{1 + \frac{2}{x} + \frac{3}{x^2}} = -x \left(1 + \frac{1}{x} + \frac{1}{x^2} + o(1/x^2)\right) = -x - 1 - \frac{1}{x} + o(1/x).
\]
The asymptote at $-\infty$ is $(\Delta'): y = -x - 1$.
The difference is $f(x) - (-x - 1) = -\frac{1}{x} + o(1/x)$.
Since $x < 0$, $-\frac{1}{x} > 0$, so $\mathcal{C}_f$ lies strictly \textbf{above} $(\Delta')$ at $-\infty$. \hfill $\blacksquare$
\end{solution}

% ===========================================================================
% EXERCISE 6.6
% ===========================================================================
\begin{exercise}[$\bigstar\bigstar$ Intermediate]{Limits of Riemann Sums}
Using the Riemann Sum Limit Theorem:
\[
\lim_{n \to \infty} \frac{1}{n} \sum_{k=1}^n f\left(\frac{k}{n}\right) = \int_0^1 f(x)\,\dd x,
\]
compute the exact limit of each of the following sequences:
\begin{enumerate}[leftmargin=1.5em,itemsep=2pt]
  \item $S_n = \sum_{k=1}^n \frac{n}{n^2 + k^2}$.
  \item $T_n = \sum_{k=1}^n \frac{1}{n + k}$.
  \item $P_n = \frac{1}{n} \sum_{k=1}^n \sqrt{1 - \left(\frac{k}{n}\right)^2}$.
\end{enumerate}
\end{exercise}

\begin{tabletworkspace}[5.0cm]
\end{tabletworkspace}

\begin{solution}
\textbf{1. Limit of $S_n$:}
Factor $\frac{1}{n}$ from the sum:
\[
S_n = \sum_{k=1}^n \frac{n}{n^2(1 + (k/n)^2)} = \frac{1}{n} \sum_{k=1}^n \frac{1}{1 + (k/n)^2}.
\]
This is a standard Riemann sum for $f(x) = \frac{1}{1 + x^2}$ on $[0, 1]$.
Since $f$ is continuous on $[0, 1]$:
\[
\lim_{n \to \infty} S_n = \int_0^1 \frac{\dd x}{1 + x^2} = [\arctan(x)]_0^1 = \arctan(1) - \arctan(0) = \frac{\pi}{4}. \quad \blacksquare
\]

\textbf{2. Limit of $T_n$:}
Factor $\frac{1}{n}$:
\[
T_n = \sum_{k=1}^n \frac{1}{n(1 + k/n)} = \frac{1}{n} \sum_{k=1}^n \frac{1}{1 + k/n}.
\]
This is a Riemann sum for $f(x) = \frac{1}{1 + x}$ on $[0, 1]$.
\[
\lim_{n \to \infty} T_n = \int_0^1 \frac{\dd x}{1 + x} = [\ln(1 + x)]_0^1 = \ln(2) - \ln(1) = \ln(2). \quad \blacksquare
\]

\textbf{3. Limit of $P_n$:}
$P_n = \frac{1}{n} \sum_{k=1}^n \sqrt{1 - (k/n)^2}$ is a Riemann sum for $f(x) = \sqrt{1 - x^2}$ on $[0, 1]$.
\[
\lim_{n \to \infty} P_n = \int_0^1 \sqrt{1 - x^2}\,\dd x.
\]
Geometrically, this integral represents the area of a quarter of the unit circle of radius 1:
\[
\int_0^1 \sqrt{1 - x^2}\,\dd x = \frac{\pi (1)^2}{4} = \frac{\pi}{4}. \quad \blacksquare
\]
\end{solution}

% ===========================================================================
% EXERCISE 6.7
% ===========================================================================
\begin{exercise}[$\bigstar\bigstar$ Intermediate]{Differentiating Integrals with Variable Bounds (Leibniz Integral Rule)}
Let $u, v: I \to \R$ be differentiable functions and $f: \R \to \R$ continuous.
Define:
\[
F(x) = \int_{u(x)}^{v(x)} f(t)\,\dd t.
\]
\begin{enumerate}[leftmargin=1.5em,itemsep=2pt]
  \item Prove using the Fundamental Theorem of Calculus and the Chain Rule that:
  \[
  F'(x) = f(v(x)) v'(x) - f(u(x)) u'(x).
  \]
  \item Calculate the derivative of $G(x) = \int_{0}^{x^2} e^{-t^2}\,\dd t$.
  \item Calculate the derivative of $H(x) = \int_{x}^{2x} \frac{\dd t}{1 + t^3}$ for $x > 0$.
\end{enumerate}
\end{exercise}

\begin{tabletworkspace}[5.0cm]
\end{tabletworkspace}

\begin{solution}
\textbf{1. Proof of Leibniz Rule:}
Let $f$ be continuous on $\R$. By the Fundamental Theorem of Calculus, $f$ admits an antiderivative $\Phi$ such that $\Phi'(t) = f(t)$.
By the evaluation formula:
\[
F(x) = \int_{u(x)}^{v(x)} f(t)\,\dd t = \Phi(v(x)) - \Phi(u(x)).
\]
Since $u$ and $v$ are differentiable and $\Phi$ is differentiable, $F$ is differentiable by the Chain Rule:
\[
F'(x) = \frac{\dd}{\dd x}[\Phi(v(x))] - \frac{\dd}{\dd x}[\Phi(u(x))] = \Phi'(v(x)) v'(x) - \Phi'(u(x)) u'(x).
\]
Since $\Phi' = f$:
\[
F'(x) = f(v(x)) v'(x) - f(u(x)) u'(x). \quad \blacksquare
\]

\textbf{2. Derivative of $G(x) = \int_0^{x^2} e^{-t^2}\,\dd t$:}
Here $u(x) = 0 \implies u'(x) = 0$, $v(x) = x^2 \implies v'(x) = 2x$, and $f(t) = e^{-t^2}$.
\[
G'(x) = e^{-(x^2)^2} \cdot (2x) - 0 = 2x e^{-x^4}.
\]

\textbf{3. Derivative of $H(x) = \int_x^{2x} \frac{\dd t}{1 + t^3}$:}
Here $u(x) = x \implies u'(x) = 1$, $v(x) = 2x \implies v'(x) = 2$, and $f(t) = \frac{1}{1 + t^3}$.
\[
H'(x) = \frac{1}{1 + (2x)^3} \cdot (2) - \frac{1}{1 + x^3} \cdot (1) = \frac{2}{1 + 8x^3} - \frac{1}{1 + x^3}. \quad \blacksquare
\]
\end{solution}

% ===========================================================================
% EXERCISE 6.8
% ===========================================================================
\begin{exercise}[$\bigstar\bigstar$ Intermediate]{Systematic Integration by Parts (IPP)}
Calculate each of the following integrals using integration by parts:
\begin{enumerate}[leftmargin=1.5em,itemsep=2pt]
  \item $I_1 = \int x^2 e^{2x}\,\dd x$.
  \item $I_2 = \int_1^e x \ln(x)\,\dd x$.
  \item $I_3 = \int e^x \cos(x)\,\dd x \quad (\text{circular / cyclic IPP})$.
\end{enumerate}
\end{exercise}

\begin{tabletworkspace}[5.0cm]
\end{tabletworkspace}

\begin{solution}
\textbf{1. Calculation of $I_1 = \int x^2 e^{2x}\,\dd x$ (Double IPP):}
\begin{itemize}[leftmargin=1.5em,itemsep=1pt]
  \item First IPP: Set $u = x^2 \implies u' = 2x$, and $v' = e^{2x} \implies v = \frac{1}{2}e^{2x}$.
  \[
  I_1 = \frac{1}{2}x^2 e^{2x} - \int x e^{2x}\,\dd x.
  \]
  \item Second IPP for $\int x e^{2x}\,\dd x$: Set $u = x \implies u' = 1$, and $v' = e^{2x} \implies v = \frac{1}{2}e^{2x}$.
  \[
  \int x e^{2x}\,\dd x = \frac{1}{2}x e^{2x} - \frac{1}{2}\int e^{2x}\,\dd x = \frac{1}{2}x e^{2x} - \frac{1}{4}e^{2x}.
  \]
  \item Combining:
  \[
  I_1 = \frac{1}{2}x^2 e^{2x} - \frac{1}{2}x e^{2x} + \frac{1}{4}e^{2x} + C = \frac{1}{4}e^{2x}(2x^2 - 2x + 1) + C. \quad \blacksquare
  \]
\end{itemize}

\textbf{2. Calculation of $I_2 = \int_1^e x \ln(x)\,\dd x$:}
Set $u = \ln(x) \implies u' = \frac{1}{x}$, and $v' = x \implies v = \frac{x^2}{2}$.
\begin{align*}
I_2 &= \left[ \frac{x^2}{2}\ln(x) \right]_1^e - \int_1^e \frac{x^2}{2} \cdot \frac{1}{x}\,\dd x = \left( \frac{e^2}{2}\ln(e) - 0 \right) - \frac{1}{2}\int_1^e x\,\dd x \\
&= \frac{e^2}{2} - \frac{1}{2}\left[ \frac{x^2}{2} \right]_1^e = \frac{e^2}{2} - \frac{1}{4}(e^2 - 1) = \frac{e^2 + 1}{4}. \quad \blacksquare
\end{align*}

\textbf{3. Calculation of $I_3 = \int e^x \cos(x)\,\dd x$ (Cyclic IPP):}
Set $u = \cos x \implies u' = -\sin x$, and $v' = e^x \implies v = e^x$.
\[
I_3 = e^x \cos x - \int e^x (-\sin x)\,\dd x = e^x \cos x + \int e^x \sin x\,\dd x.
\]
Apply IPP to $\int e^x \sin x\,\dd x$: Set $u = \sin x \implies u' = \cos x$, and $v' = e^x \implies v = e^x$.
\[
\int e^x \sin x\,\dd x = e^x \sin x - \int e^x \cos x\,\dd x = e^x \sin x - I_3.
\]
Substituting back into $I_3$:
\[
I_3 = e^x \cos x + e^x \sin x - I_3 \iff 2I_3 = e^x(\cos x + \sin x).
\]
Therefore:
\[
I_3 = \frac{e^x}{2}(\sin x + \cos x) + C. \quad \blacksquare
\]
\end{solution}

% ===========================================================================
% EXERCISE 6.9
% ===========================================================================
\begin{exercise}[$\bigstar\bigstar$ Intermediate]{Integration by Substitution (Change of Variables)}
Calculate the following integrals using an appropriate substitution:
\begin{enumerate}[leftmargin=1.5em,itemsep=2pt]
  \item $J_1 = \int_0^1 \frac{x}{\sqrt{1 + x^2}}\,\dd x \quad (\text{substitute } u = 1 + x^2)$.
  \item $J_2 = \int_0^{\pi/2} \frac{\cos(x)}{1 + \sin^2(x)}\,\dd x \quad (\text{substitute } u = \sin x)$.
  \item $J_3 = \int_1^4 \frac{\dd x}{x(1 + \sqrt{x})} \quad (\text{substitute } u = \sqrt{x})$.
\end{enumerate}
\end{exercise}

\begin{tabletworkspace}[5.0cm]
\end{tabletworkspace}

\begin{solution}
\textbf{1. Calculation of $J_1$:}
Let $u = 1 + x^2 \implies \dd u = 2x\,\dd x \iff x\,\dd x = \frac{\dd u}{2}$.
Bounds: when $x = 0$, $u = 1$; when $x = 1$, $u = 2$.
\[
J_1 = \int_1^2 \frac{\dd u / 2}{\sqrt{u}} = \frac{1}{2} \int_1^2 u^{-1/2}\,\dd u = \frac{1}{2} \left[ 2\sqrt{u} \right]_1^2 = [\sqrt{u}]_1^2 = \sqrt{2} - 1. \quad \blacksquare
\]

\textbf{2. Calculation of $J_2$:}
Let $u = \sin(x) \implies \dd u = \cos(x)\,\dd x$.
Bounds: when $x = 0$, $u = \sin 0 = 0$; when $x = \pi/2$, $u = \sin(\pi/2) = 1$.
\[
J_2 = \int_0^1 \frac{\dd u}{1 + u^2} = [\arctan(u)]_0^1 = \arctan(1) - \arctan(0) = \frac{\pi}{4}. \quad \blacksquare
\]

\textbf{3. Calculation of $J_3$:}
Let $u = \sqrt{x} \iff x = u^2 \implies \dd x = 2u\,\dd u$.
Bounds: when $x = 1$, $u = 1$; when $x = 4$, $u = 2$.
Substitute:
\[
J_3 = \int_1^2 \frac{2u\,\dd u}{u^2(1 + u)} = 2 \int_1^2 \frac{\dd u}{u(1 + u)}.
\]
Decompose into partial fractions:
\[
\frac{1}{u(1 + u)} = \frac{1}{u} - \frac{1}{1 + u}.
\]
Integrate:
\begin{align*}
J_3 &= 2 \int_1^2 \left( \frac{1}{u} - \frac{1}{1 + u} \right)\dd u = 2 [\ln|u| - \ln|1 + u|]_1^2 = 2 \left[ \ln\left(\frac{u}{1 + u}\right) \right]_1^2 \\
&= 2 \left( \ln(2/3) - \ln(1/2) \right) = 2 \ln\left(\frac{2/3}{1/2}\right) = 2 \ln\left(\frac{4}{3}\right). \quad \blacksquare
\end{align*}
\end{solution}

% ===========================================================================
% EXERCISE 6.10
% ===========================================================================
\begin{exercise}[$\bigstar\bigstar\bigstar$ Advanced / Exam Level]{Partial Fraction Decomposition of Rational Functions}
Consider the rational fraction:
\[
f(x) = \frac{2x^2 - 3x + 5}{(x - 1)(x^2 + 1)}.
\]
\begin{enumerate}[leftmargin=1.5em,itemsep=2pt]
  \item Determine the partial fraction decomposition of $f(x)$ in $\R(X)$:
  \[
  \frac{2x^2 - 3x + 5}{(x - 1)(x^2 + 1)} = \frac{A}{x - 1} + \frac{Bx + C}{x^2 + 1}.
  \]
  \item Determine constants $A, B, C$ by identifying coefficients or evaluating at special values.
  \item Calculate the indefinite integral:
  \[
  I = \int \frac{2x^2 - 3x + 5}{(x - 1)(x^2 + 1)}\,\dd x.
  \]
\end{enumerate}
\end{exercise}

\begin{tabletworkspace}[5.0cm]
\end{tabletworkspace}

\begin{solution}
\textbf{1. Decomposition Form:}
The denominator has a real linear factor $(x - 1)$ and an irreducible quadratic factor $(x^2 + 1)$ with $\Delta = -4 < 0$. The degree of the numerator (2) is strictly less than the degree of the denominator (3).
The decomposition form over $\R$ is:
\[
\frac{2x^2 - 3x + 5}{(x - 1)(x^2 + 1)} = \frac{A}{x - 1} + \frac{Bx + C}{x^2 + 1}.
\]

\textbf{2. Determination of Coefficients $A, B, C$:}
\begin{itemize}[leftmargin=1.5em,itemsep=2pt]
  \item \textbf{Finding $A$:} Multiply both sides by $(x - 1)$ and set $x = 1$:
  \[
  A = \left. \frac{2x^2 - 3x + 5}{x^2 + 1} \right|_{x=1} = \frac{2(1)^2 - 3(1) + 5}{1^2 + 1} = \frac{4}{2} = 2.
  \]
  \item \textbf{Finding $C$:} Set $x = 0$ in the identity:
  \[
  \frac{5}{(-1)(1)} = \frac{A}{-1} + \frac{C}{1} \implies -5 = -2 + C \implies C = -3.
  \]
  \item \textbf{Finding $B$:} Multiply both sides by $x$ and take the limit as $x \to \infty$:
  \[
  \lim_{x \to \infty} \frac{2x^3 - 3x^2 + 5x}{(x - 1)(x^2 + 1)} = 2 = A + B \implies 2 = 2 + B \implies B = 0.
  \]
\end{itemize}
Thus $A = 2$, $B = 0$, $C = -3$.
The partial fraction decomposition is:
\[
\frac{2x^2 - 3x + 5}{(x - 1)(x^2 + 1)} = \frac{2}{x - 1} - \frac{3}{x^2 + 1}.
\]

\textbf{3. Integration:}
\begin{align*}
I &= \int \left( \frac{2}{x - 1} - \frac{3}{x^2 + 1} \right)\dd x \\
&= 2 \int \frac{\dd x}{x - 1} - 3 \int \frac{\dd x}{x^2 + 1} \\
&= 2 \ln|x - 1| - 3 \arctan(x) + C. \quad \blacksquare
\end{align*}
\end{solution}

% ===========================================================================
% EXERCISE 6.11
% ===========================================================================
\begin{exercise}[$\bigstar\bigstar\bigstar$ Advanced / Exam Level]{The Wallis Integrals and Asymptotic Formula for $\pi$}
Define the sequence of Wallis integrals for $n \in \N$ by:
\[
W_n = \int_0^{\pi/2} \sin^n(x)\,\dd x.
\]
\begin{enumerate}[leftmargin=1.5em,itemsep=2pt]
  \item Compute $W_0$ and $W_1$.
  \item Using integration by parts on $\sin^n(x) = \sin(x) \sin^{n-1}(x)$, establish the recurrence relation:
  \[
  W_n = \frac{n - 1}{n} W_{n-2} \quad \text{for all } n \ge 2.
  \]
  \item Deduce explicit formulas for $W_{2p}$ and $W_{2p+1}$ in terms of double factorials:
  \[
  W_{2p} = \frac{(2p)!}{2^{2p} (p!)^2} \frac{\pi}{2}, \qquad W_{2p+1} = \frac{2^{2p} (p!)^2}{(2p + 1)!}.
  \]
  \item Prove that $(n + 1) W_{n+1} W_n = \frac{\pi}{2}$ for all $n \ge 0$, and deduce the Wallis product formula for $\pi$:
  \[
  \lim_{p \to \infty} \frac{1}{2p} \left( \frac{2 \cdot 4 \cdot 6 \dots (2p)}{1 \cdot 3 \cdot 5 \dots (2p - 1)} \right)^2 = \frac{\pi}{2}.
  \]
\end{enumerate}
\end{exercise}

\begin{tabletworkspace}[5.5cm]
\end{tabletworkspace}

\begin{solution}
\textbf{1. Initial Terms:}
\[
W_0 = \int_0^{\pi/2} 1\,\dd x = \frac{\pi}{2}, \qquad W_1 = \int_0^{\pi/2} \sin(x)\,\dd x = [-\cos x]_0^{\pi/2} = 0 - (-1) = 1.
\]

\textbf{2. Recurrence Relation via IPP:}
For $n \ge 2$, write $W_n = \int_0^{\pi/2} \sin(x) \sin^{n-1}(x)\,\dd x$.
Set $u = \sin^{n-1}(x) \implies u' = (n - 1)\cos(x)\sin^{n-2}(x)$, and $v' = \sin(x) \implies v = -\cos(x)$.
\begin{align*}
W_n &= [-\cos(x)\sin^{n-1}(x)]_0^{\pi/2} + (n - 1)\int_0^{\pi/2} \cos^2(x)\sin^{n-2}(x)\,\dd x \\
&= 0 + (n - 1)\int_0^{\pi/2} (1 - \sin^2 x)\sin^{n-2}(x)\,\dd x \\
&= (n - 1) W_{n-2} - (n - 1) W_n.
\end{align*}
Rearranging:
\[
W_n + (n - 1)W_n = (n - 1)W_{n-2} \iff n W_n = (n - 1)W_{n-2} \iff W_n = \frac{n - 1}{n} W_{n-2}.
\]

\textbf{3. Closed-Form Expressions:}
\begin{itemize}[leftmargin=1.5em,itemsep=2pt]
  \item For $n = 2p$:
  \[
  W_{2p} = \frac{2p - 1}{2p} \cdot \frac{2p - 3}{2p - 2} \dots \frac{1}{2} W_0 = \frac{(2p - 1)!!}{(2p)!!} \frac{\pi}{2} = \frac{(2p)!}{2^{2p} (p!)^2} \frac{\pi}{2}.
  \]
  \item For $n = 2p + 1$:
  \[
  W_{2p+1} = \frac{2p}{2p + 1} \cdot \frac{2p - 2}{2p - 1} \dots \frac{2}{3} W_1 = \frac{(2p)!!}{(2p + 1)!!} = \frac{2^{2p} (p!)^2}{(2p + 1)!}.
  \]
\end{itemize}

\textbf{4. Product Identity and Wallis Formula:}
Compute the product $(n + 1) W_{n+1} W_n$:
For $n = 2p$:
\[
(2p + 1) W_{2p+1} W_{2p} = (2p + 1) \cdot \frac{2^{2p}(p!)^2}{(2p + 1)!} \cdot \frac{(2p)!}{2^{2p}(p!)^2}\frac{\pi}{2} = \frac{\pi}{2}.
\]
Since $0 < \sin x < 1$ on $(0, \pi/2)$, $\sin^{n+1} x < \sin^n x < \sin^{n-1} x$, which implies $W_{n+1} < W_n < W_{n-1}$.
Dividing by $W_{n+1}$:
\[
1 < \frac{W_n}{W_{n+1}} < \frac{W_{n-1}}{W_{n+1}} = \frac{n + 1}{n} \xrightarrow{n \to \infty} 1.
\]
By the Squeeze Theorem, $W_n \sim W_{n+1}$.
Thus $\frac{\pi}{2} = (n + 1)W_{n+1} W_n \sim n W_n^2 \implies W_n \sim \sqrt{\frac{\pi}{2n}}$.
Substituting $n = 2p$:
\[
\frac{W_{2p}}{W_{2p+1}} = \left( \frac{(2p)!!}{(2p - 1)!!} \right)^2 \frac{1}{2p + 1}\frac{\pi}{2} \xrightarrow{p \to \infty} 1.
\]
Rearranging yields the historic Wallis product formula for $\pi$:
\[
\frac{\pi}{2} = \lim_{p \to \infty} \prod_{k=1}^p \frac{4k^2}{4k^2 - 1} = \lim_{p \to \infty} \left( \frac{2 \cdot 4 \cdot 6 \dots (2p)}{1 \cdot 3 \cdot 5 \dots (2p - 1)} \right)^2 \frac{1}{2p + 1}. \quad \blacksquare
\]
\end{solution}

% ===========================================================================
% EXERCISE 6.12
% ===========================================================================
\begin{exercise}[$\bigstar\bigstar\bigstar$ Advanced / Exam Level]{Synthesis Problem: Error Bounds in Numerical Quadrature (Trapezoidal Rule)}
Let $f \in C^2([a, b])$ with $M_2 = \max_{x \in [a, b]} |f''(x)|$.
Approximate the integral $I = \int_a^b f(x)\,\dd x$ using the composite Trapezoidal Rule on $n$ subintervals of length $h = \frac{b - a}{n}$:
\[
T_n = \frac{h}{2} \left[ f(a) + 2\sum_{k=1}^{n-1} f(a + kh) + f(b) \right].
\]
\begin{enumerate}[leftmargin=1.5em,itemsep=2pt]
  \item On an elementary interval $[x_k, x_{k+1}]$ with midpoint $m_k = \frac{x_k + x_{k+1}}{2}$, express $f(x)$ via Taylor's formula around $m_k$.
  \item Prove that the elementary error on $[x_k, x_{k+1}]$ satisfies:
  \[
  \left| \int_{x_k}^{x_{k+1}} f(x)\,\dd x - \frac{h}{2}[f(x_k) + f(x_{k+1})] \right| \le \frac{M_2 h^3}{12}.
  \]
  \item Sum over all $n$ subintervals to establish the global error bound:
  \[
  |I - T_n| \le \frac{M_2 (b - a)^3}{12 n^2}.
  \]
  \item Application: How many subintervals $n$ are required to approximate $\int_0^1 e^{-x^2}\,\dd x$ using the Trapezoidal Rule with an absolute error strictly less than $10^{-4}$?
\end{enumerate}
\end{exercise}

\begin{tabletworkspace}[5.5cm]
\end{tabletworkspace}

\begin{solution}
\textbf{1. Taylor Expansion around Midpoint $m_k$:}
Let $h = x_{k+1} - x_k$. The midpoint is $m_k = x_k + h/2$.
For $x \in [x_k, x_{k+1}]$, let $t = x - m_k \in [-h/2, h/2]$.
By Taylor's formula with integral remainder:
\[
f(m_k + t) = f(m_k) + t f'(m_k) + \int_0^t (t - s) f''(m_k + s)\,\dd s.
\]

\textbf{2. Elementary Trapezoidal Error:}
Integrating $f(m_k + t)$ over $t \in [-h/2, h/2]$:
\[
\int_{x_k}^{x_{k+1}} f(x)\,\dd x = \int_{-h/2}^{h/2} f(m_k + t)\,\dd t = h f(m_k) + 0 + \int_{-h/2}^{h/2} \left( \int_0^t (t - s)f''(m_k + s)\,\dd s \right)\dd t.
\]
The trapezoidal approximation is:
\[
\frac{h}{2}[f(x_k) + f(x_{k+1})] = \frac{h}{2}\left[ f(m_k - h/2) + f(m_k + h/2) \right] = h f(m_k) + \frac{h^3}{8} f''(m_k) + O(h^5).
\]
Subtracting and bounding with $|f''(x)| \le M_2$:
Standard integration yields:
\[
\left| \int_{x_k}^{x_{k+1}} f(x)\,\dd x - \frac{h}{2}[f(x_k) + f(x_{k+1})] \right| \le \frac{M_2 h^3}{12}.
\]

\textbf{3. Global Error Bound:}
Summing over all $n$ subintervals $k = 0, 1, \dots, n - 1$:
\[
|I - T_n| \le \sum_{k=0}^{n-1} \frac{M_2 h^3}{12} = n \cdot \frac{M_2 h^3}{12}.
\]
Substitute $h = \frac{b - a}{n}$:
\[
|I - T_n| \le n \cdot \frac{M_2}{12}\left(\frac{b - a}{n}\right)^3 = \frac{M_2 (b - a)^3}{12 n^2}. \quad \blacksquare
\]

\textbf{4. Application to $\int_0^1 e^{-x^2}\,\dd x$:}
Here $[a, b] = [0, 1]$, so $b - a = 1$.
Compute derivatives of $f(x) = e^{-x^2}$:
\[
f'(x) = -2x e^{-x^2}, \qquad f''(x) = (4x^2 - 2)e^{-x^2}.
\]
On $[0, 1]$:
\begin{itemize}[leftmargin=1.5em,itemsep=1pt]
  \item At $x = 0$: $|f''(0)| = |-2| = 2$.
  \item At $x = 1$: $|f''(1)| = |2 e^{-1}| = \frac{2}{e} \approx 0.736$.
  \item Derivative of $f''$: $f'''(x) = (12x - 8x^3)e^{-x^2} = 4x(3 - 2x^2)e^{-x^2} = 0 \implies x = \sqrt{3/2} \approx 1.22 > 1$.
\end{itemize}
Thus the maximum of $|f''(x)|$ on $[0, 1]$ occurs at $x = 0$:
\[
M_2 = \max_{x \in [0, 1]} |f''(x)| = 2.
\]
The error bound gives:
\[
|I - T_n| \le \frac{2 \times (1)^3}{12 n^2} = \frac{1}{6n^2}.
\]
We require:
\[
\frac{1}{6n^2} < 10^{-4} \iff 6n^2 > 10^4 = 10000 \iff n^2 > \frac{10000}{6} \approx 1666.67.
\]
Taking square roots:
\[
n > \sqrt{1666.67} \approx 40.82.
\]
Therefore, choosing $n = 41$ subintervals guarantees an error strictly less than $10^{-4} = 0.0001$. \hfill $\blacksquare$
\end{solution}
